% ===== §5 =====
\section{The Aesthetic Measure of the Ethical Response}
\label{p26:sec:ethics}

\noindent
This section takes the fourth dimension, the ethical demand a relation makes, and establishes that at its root the ethical response rests on an aesthetic measure, a sensed fittingness that cannot be turned into a proposition, so that ethics is a matter of cultivated perception rather than of applied doctrine. The section states the phenomenon, reviews the singularity of the other and the aesthetic register through which it is met, isolates the core, and draws the claim. This is the dimension the paper marks as ethical at its root, and the review is correspondingly fuller.

\subsection*{The phenomenon}

A relation makes ethical demands, and the first of them is a demand to answer this person, the one before me, in their singularity. Language seems able to meet the demand. We have the resources of moral speech: principles, reasons, the vocabulary of obligation and desert. It appears that to act well toward another is to apply the right rule to the case, and that the rule and the case can both be said. A long tradition in moral philosophy proceeds on exactly this appearance, treating the ethical task as the subsumption of a particular case under a general principle.

They can be said only up to a point, and the point is where ethics begins.

\subsection*{The singularity of the other}

To grasp the other through any predicate, any category, any rule, is already to grasp them as an instance of a kind, a bearer of the stated features, one who falls under the description. Levinas made this the center of ethics: the other faces me not as an instance of humanity in general but as this one, whose claim on me precedes and exceeds every category under which I might bring them, so that the ethical relation is a relation to what no concept contains \parencite{levinas1969totality}. The face, in his vocabulary, is exactly what resists becoming a theme, a content I could state and thereby master. The ethical claim comes from what the predicates do not exhaust, from this other as this other and not as a case of a type. So far the limit is a limit on conceptual grasp: the singular exceeds the general, and the other's claim issues from the excess.

But this leaves a question Levinas does not fully answer, and it is the question that carries the section. If the other exceeds every rule, how is the fitting response to them found? It cannot be deduced, for deduction works from generals and the case is singular. What faculty, then, takes the measure of a singular to which no rule reaches?

\subsection*{The aesthetic register}

The answer this paper proposes is that the measure is taken aesthetically, in the register of \emph{aisthesis}, sensed before it is stated and not reducible to any rule from which it could be read off. How much to say and how much to leave unsaid, when to press and when to hold back, the exact weight of a gesture, the timing of a silence: all of this is apprehended in the way one apprehends the rightness of a line or the wrongness of a false note. It is a perception of fittingness, not an inference to it.

Two sources make this register precise. Merleau-Ponty showed that perception is not the passive reception of data but an active taking-up of a situation, a bodily attunement that grasps a whole before it articulates any part, and that this attunement has its own intelligence, a knowing in the perceiving that no subsequent proposition exhausts \parencite{merleauponty1962phenomenology}. The fitting response to another is grasped by exactly this perceptual intelligence: one senses the shape of the situation and what it calls for, as a whole, before and beneath any rule one could state. Dufrenne, extending phenomenology to aesthetic experience, described the perception of an aesthetic object as the apprehension of an expressed meaning that is given only in the sensing and cannot be detached from it into a concept, a meaning that is felt as the depth of the perceived and lost the moment one tries to state it apart from the perceiving \parencite{dufrenne1973phenomenology}. The measure of the ethical response has this form. It is an expressed meaning of the singular situation, given in the perception of it and not detachable into a rule.

This is why the ethical response cannot be propositionalized. The proposition states a general, and the fittingness at stake is of this response, here, to this one. To translate the sensed measure into a statable rule is to replace the perception of a singular with the application of a general, which is to pass over exactly the singularity the other's claim consisted in. The measure lives in the sensing and dies in the saying.

\subsection*{The core}

The core language cannot reach here is the \emph{aesthetic measure} of the ethical response, the sensed fittingness that cannot be turned into a proposition. It is not that the right rule is hard to find or hard to phrase. It is that the fittingness at stake is of a kind with the fittingness of a form, apprehended in perception and lost in translation to a rule, because the rule states a general and the fittingness is singular. The ethical, at its root, is answerable to a measure that is aesthetic in the strict sense: taken by \emph{aisthesis}, in the perceiving, and not by the understanding, in the proposition.

\begin{claim}[The aesthetic root of the ethical]
The fitting response to the singular other cannot be derived from any rule, because rules state generals and the fittingness is of this response to this one. It is grasped aesthetically, in the register of aisthesis, as one grasps the rightness of a form, and is lost when translated into a proposition. Ethics at its root rests on a perception, not a doctrine.
\end{claim}

\begin{casebox}{the timing no rule gives}
Someone before you is in pain, and the question is whether to speak or to keep silent, and if to speak, in what weight and at what moment. No rule delivers the answer, for the rule states a general and this is this person, now. What guides the fitting response is a perception of the whole situation, a sensed measure taken in the way one senses that a further word would be one word too many. Those who respond well here are not those who know more rules but those whose perception has been schooled finer, and what they have is not a doctrine but an eye.
\end{casebox}

\subsection*{Why the limit is generative here}

This is what keeps ethics a matter of cultivated perception rather than of applied doctrine. Because the fitting response cannot be propositionalized, it must be perceived, and the ethical life is the schooling of a perception fine enough to take the measure of a singular case. Consider the counter-case, a morality fully written as rules. One would discharge it by looking up the rule and applying it, and the looking-up would pass over the singularity of the other in favor of the type under which the rule brought them; a perfectly codified ethics would be an ethics that never met anyone, only cases. The limit forbids this. Because the measure is aesthetic and cannot be codified, ethics is returned to the trained eye and the trained ear, to a sense of the fitting that no code replaces, and the ethical life becomes the long cultivation of that sense. This connects the ethical dimension to the aesthetic education the series has developed elsewhere, and it will return in Part Four, where the ethics of making room for the other's unsymbolizable core is shown to rest on exactly this perceptual measure. What cannot be said in the ethical relation is not a failing of moral language. It is the opening in which answering this other, and not merely applying a rule to a case, becomes possible at all.
